\section{Panduan Sintaks \& Fitur Utama}
\label{sec:sintaks}

\subsection{Penulisan Rumus Matematika}
LaTeX memiliki kemampuan luar biasa dalam merender rumus matematika:
\begin{equation}
    f(x) = \int_{-\infty}^{\infty} \frac{1}{\sqrt{2\pi\sigma^2}} e^{-\frac{(x-\mu)^2}{2\sigma^2}} dx = 1
    \label{eq:normal_dist}
\end{equation}
Persamaan di atas dapat dirujuk secara otomatis menggunakan label \eqref{eq:normal_dist}.

\subsection{Manajemen Sitasi Modern dengan Biber}
Dengan \texttt{biblatex}, sitasi dapat dipanggil menggunakan perintah \texttt{\textbackslash cite}. Contohnya, karya fundamental Donald Knuth mengenai TeX \cite{knuth1984texbook} dan Leslie Lamport mengenai LaTeX \cite{lamport1994latex}.

\subsection{Struktur Tabel}
Tabel rapi dapat dibuat dengan paket \texttt{booktabs}:
\begin{table}[htbp]
\centering
\caption{Perbandingan Format Dokumen Tradisional vs Modular}
\label{tab:perbandingan}
\begin{tabular}{lll}
\toprule
\textbf{Parameter} & \textbf{Satu File (Flat)} & \textbf{Modular} \\
\midrule
Navigasi Bab & Sulit & Sangat Mudah \\
Kolaborasi Git & Rawan Konflik & Aman \& Terisolasi \\
Kompilasi & Lambat & Cepat via Root File \\
\bottomrule
\end{tabular}
\end{table}
